\section{Coherence and loss of quantum information}
\label{sec:platform-coherence}

Keeping an electron in its dot does not by itself preserve a spin
superposition. Its relative phase depends on the energy splitting
throughout the wait or gate. If that splitting fluctuates in a way
we do not know, repetitions of the same control sequence need not
produce the same phase. Other processes change spin populations or
populate states outside the chosen qubit space. These mechanisms
have different consequences for the operations derived above.

\subsection{Phase accumulation, dephasing, and spin echo}

Consider a single-spin qubit with the drive switched off. In a frame
rotating at a fixed reference frequency, write the longitudinal
Hamiltonian in our logical basis as
\begin{equation}
 H_{\mathrm{wait}}(t)=-\frac{\hbar}{2}
                 [\delta_0+\xi(t)]Z,
 \qquad
 \phi(t)=\int_0^t\xi(t')\,dt'.
 \label{eq:platform-wait-noise}
\end{equation}
Here $\delta_0$ is the mean spin frequency minus the reference
frequency, and $\xi(t)$ is a classical fluctuation of that difference,
with units of angular frequency. Magnetic fluctuations can change
the local field; electric fluctuations can change the effective
$g$ factor or the position of the electron in a field gradient.
The model retains only the resulting longitudinal frequency noise.
It does not include spin flips or excitation of other electronic
states.

For a given realization of the noise, the two logical amplitudes
acquire opposite phases. Consequently the density-matrix element
$\rho_{01}=\bra0\rho\ket1$ is multiplied by
$e^{i\delta_0t+i\phi(t)}$, while both populations stay unchanged.
Averaging over repetitions gives
\begin{equation}
 \overline{\rho}_{01}(t)
       =e^{i\delta_0t}W(t)\rho_{01}(0),
 \qquad W(t)=\left\langle e^{i\phi(t)}\right\rangle.
 \label{eq:platform-phase-coherence}
\end{equation}
The brackets denote an average over noise realizations sampled in
the experiment. The measured mean detuning determines the correction
to the reference frequency, or the phase tracked during the wait
(Appendix~\ref{sec:idle-single-spin-characterization}).
A spread of unknown phases reduces $|W|$ and
hence the coherence of the averaged state. This is dephasing, even
though every realization of this particular model evolves unitarily.

The analysis rotation in
Eq.~\eqref{eq:platform-analysis-probability} makes this reduction
observable. Prepare $\ket+$, wait for time $t$, apply $R_y(-\pi/2)$,
and measure in the logical basis. Since
$\operatorname{Tr}(\rho X)=2\operatorname{Re}\rho_{01}$, the zero
probability is
\begin{equation}
 p_0(t)=\frac{1+\operatorname{Re}[e^{i\delta_0t}W(t)]}{2}.
 \label{eq:platform-ramsey-coherence}
\end{equation}
This is a Ramsey measurement. The known detuning produces fringes;
phase fluctuations reduce their contrast. In a simple limit,
$\xi(t)=\xi$ is constant during each wait but varies between
repetitions with a zero-mean Gaussian distribution of variance
$\sigma_\omega^2$. The coherence and measured probability are then
\begin{equation}
 \begin{aligned}
 W(t)&=e^{-\sigma_\omega^2t^2/2},\\
 p_0(t)&=\frac{1+e^{-(t/T_2^*)^2}\cos(\delta_0t)}2,
 \qquad T_2^*=\frac{\sqrt2}{\sigma_\omega}.
 \end{aligned}
 \label{eq:platform-gaussian-coherence}
\end{equation}
Thus $T_2^*$ is the $1/e$ time of the free-evolution contrast in
this model. Other noise statistics give other envelopes, and noise
slower than the duration of the data acquisition may appear as a
fixed offset instead of a sampled distribution.

A midpoint $\pi$ rotation about $X$ exchanges $\ket0$ and $\ket1$.
Its frequency, phase, and pulse area use the calibrated transverse
controls in Section~\ref{sec:single-qubit-gates}.
A component that accumulated one sign of phase in the first half
therefore accumulates the opposite sign in the second half.
In a frame that tracks this known pulse, the random phase becomes
\begin{equation}
 \phi_{\mathrm e}(t)
   =\int_0^{t/2}\xi(t')\,dt'
       -\int_{t/2}^{t}\xi(t')\,dt',
 \qquad W_{\mathrm e}(t)=\left\langle e^{i\phi_{\mathrm e}(t)}\right\rangle.
 \label{eq:platform-echo-coherence}
\end{equation}
A constant offset contributes equally to the two integrals and
cancels, including $\delta_0$. Slowly varying noise is partly
cancelled; noise that changes substantially between the two halves
need not be. This spin-echo sequence can therefore preserve contrast
longer than free evolution. Its measured decay time depends on the
pulse sequence and the noise, so it must be distinguished from
$T_2^*$. Appendix~\ref{sec:measurement-ramsey-echo} develops the pulse
phases and measured probabilities for Ramsey and echo experiments.

\subsection{Relaxation and leakage}

Longitudinal phase noise is only one way of losing quantum
information. Coupling to the environment can transfer energy and
change the populations of $\ket0$ and $\ket1$. The relaxation law
in Eq.~\eqref{eq:thermal-spin-initialization} describes their
approach to thermal equilibrium on the timescale $T_1$. The same
process that helps prepare the ground spin state destroys an
arbitrary stored superposition. Exchanging the two states halfway
through a wait does not reverse an uncontrolled energy transfer
to the environment, so spin echo does not in general remove this
relaxation.

A separate possibility is population outside the computational
space. For the projector $P_{\mathrm q}$ defined in the introduction
to Section~\ref{sec:hamiltonians-to-gates}, the population outside
that space is
\begin{equation}
 L(t)=1-\operatorname{Tr}[P_{\mathrm q}\rho(t)],
 \qquad \operatorname{Tr}\rho(t)=1.
 \label{eq:platform-leakage-population}
\end{equation}
Examples include unwanted orbital or valley excitation, or
occupation of the $S=3/2$ sector of a three-spin exchange-only
qubit. Pure dephasing and relaxation between the two retained
logical states both have $L=0$. Conversely, a leaked state may
produce the same charge-readout outcome as a logical state, as the
pair-singlet test above illustrates. Leakage must therefore be
specified relative to the chosen computational space and cannot
always be inferred from a single readout population.

These distinctions also explain why one coherence time is not a
complete measure of gate quality. A gate can have a reproducible
angle error without appreciable idle dephasing, or drive leakage
even when the idle spin is long lived. The relevant error depends
on the state, control pulse, and observation being considered.
For exchange control, Eq.~\eqref{eq:exchange-detuning} gives one
useful consequence: at $\varepsilon=0$, the first derivative
$\partial J/\partial\varepsilon$ vanishes when $t_c,U,V$ are fixed.
Small detuning fluctuations then enter only at higher order, while
barrier fluctuations still change $t_c$ and hence the rotation angle.
Locate the low-sensitivity setting from the measured exchange
response along the compensated detuning-voltage path, then tune
the barrier to obtain the required exchange strength
(Appendix~\ref{app:noise-leakage}).
Avoiding leakage additionally requires limiting unwanted transitions
across charge, orbital, valley, or encoding gaps during the pulse.
Appendix~\ref{app:noise-leakage} quantifies the angle sensitivity in
Eq.~\eqref{eq:calibration-angle-noise} and the unwanted transition
amplitudes in Eq.~\eqref{eq:calibration-leakage-amplitude}.

